\documentclass[11pt]{article}

\usepackage{amsmath, amssymb}
\usepackage{geometry}
\usepackage{booktabs}
\usepackage{hyperref}
\usepackage{parskip}
\usepackage{titlesec}
\usepackage{array}
\usepackage{xcolor}

\geometry{margin=1.25in}

\title{Circuit Analysis: Return-to-Rest Constraints\\
\large C.\ elegans Locomotion Network}
\author{Distribution of Oscillators Project}
\date{\today}

\begin{document}

\maketitle

\begin{abstract}
The piecewise-linear Yuval neuron model is intrinsically bistable, meaning that once a neuron
activates it has no endogenous mechanism to return to rest.  In a half-center oscillator this is
resolved by mutual inhibition; in the full 102-cell \textit{C.~elegans} locomotion network the
situation is more complex.  This document analyses the network connectivity to identify which cell
classes lack direct inhibitory input, characterises the indirect return paths available to them,
and derives the resulting constraints on the synaptic conductance parameters.
\end{abstract}

% ──────────────────────────────────────────────────────────────────────────────
\section{Bistability of the Yuval Neuron}

The membrane voltage of each cell obeys
\begin{equation}
    C\,\dot{V} = g\,f(V) + I_{\mathrm{app}} + I_{\mathrm{syn}},
\end{equation}
where $f(V)$ is the piecewise-linear nonlinearity
\begin{equation}
    f(V) =
    \begin{cases}
        m_1(-V + L)             & V < L,\\
        m_2(V - L)(V - T)       & L \le V < T,\\
        m_3(V - H)(V - T)       & T \le V < H,\\
        m_4(-V + H)             & V \ge H,
    \end{cases}
\end{equation}
with $L = -70\ \mathrm{mV}$ (resting potential), $T = -45\ \mathrm{mV}$ (threshold),
$H = -35\ \mathrm{mV}$ (plateau potential), and slopes $m_1 = 0.7$, $m_2 = 1/81$,
$m_3 = -1/30$, $m_4 = 0.17$.

Setting $I_{\mathrm{syn}} = 0$, the fixed points satisfy $g\,f(V^*) + I_{\mathrm{app}} = 0$.
The function $f$ has zeros at $L$, $T$, and $H$; a stability analysis of the sign of $f$ in each
interval reveals two \emph{stable} fixed points ($L$ and $H$) and one \emph{unstable} fixed point
($T$):

\begin{center}
\begin{tabular}{clll}
\toprule
Region & $f(V)$ sign & Voltage trend & Fixed-point character \\
\midrule
$V < L$       & $+$ & $\uparrow$ toward $L$ & \\
$L < V < T$   & $-$ & $\downarrow$ toward $L$ & $L$: stable \\
$T < V < H$   & $+$ & $\uparrow$ toward $H$ & $T$: unstable \\
$V > H$       & $-$ & $\downarrow$ toward $H$ & $H$: stable \\
\bottomrule
\end{tabular}
\end{center}

\subsection*{Monostability threshold}

Bistability is lost when $I_{\mathrm{app}}$ is large enough to eliminate the resting fixed point.
The minimum of $g\,f(V)$ on $[L, T]$ is achieved at the midpoint and equals
\begin{equation}
    \min_{V\in[L,T]} g\,f(V)
    = g \cdot m_2 \cdot \frac{(T-L)^2}{4} \cdot (-1)
    \approx -1.93\ \mathrm{nS\cdot mV}.
\end{equation}
When $I_{\mathrm{app}} > 1.93$, the horizontal line $-I_{\mathrm{app}}/g$ lies below this minimum
and the resting fixed point disappears: the neuron is \emph{monostable at the plateau}.  This
distinction is critical for inhibitory interneurons (see Section~\ref{sec:return}).

% ──────────────────────────────────────────────────────────────────────────────
\section{Return-to-Rest in Isolated Circuits}

\subsection{Half-center oscillator}

In a two-cell mutual-inhibition circuit the return mechanism is well-defined.  Labelling the
currently active cell~0 and the suppressed cell~1:

\begin{enumerate}
    \item Cell~0 is at plateau $H$; its synaptic fatigue variable depletes,
          $\dot{s}_0 = -a(V_0 - L) < 0$.
    \item As $s_0 \to 0$ the inhibitory current onto cell~1 weakens; cell~1, driven by
          $I_{\mathrm{app}} > 0$, crosses $T$ and activates.
    \item Cell~1's fatigue $s_1$ is fully recovered (it was at rest, so
          $\dot{s}_1 = b > 0$); it delivers a strong inhibitory kick to cell~0.
    \item That kick drives $V_0$ below $T$, releasing cell~0 to rest.
\end{enumerate}

This is a \emph{release} mechanism (Wang \& Rinzel, 1992): the active cell releases the suppressed
cell by depleting its own synapse, not by the suppressed cell escaping through the inhibition.

\subsection{Excitatory–inhibitory (E/I) loop}

Consider the basic motif $E \xrightarrow{\mathrm{exc}} I \xrightarrow{\mathrm{inh}} E$.
Cell~E activates, drives cell~I above threshold, and cell~I provides the inhibitory kick that
returns cell~E to rest.

Cell~I is now at the plateau.  Whether it returns to rest depends on its tonic drive:
\begin{itemize}
    \item \textbf{$I$ is monostable at rest} ($I_{\mathrm{app},I}$ low): once the excitatory
          drive from E fatigues ($s_E \to 0$), cell~I loses its only input and passively returns
          to $L$.  No inhibitory kick is needed.
    \item \textbf{$I$ is bistable}: cell~I latches at $H$ regardless of whether its excitatory
          input disappears.  A second inhibitory synapse onto~I is required to kick it back.
\end{itemize}

The OscillatorComb files encode exactly this distinction: cells designated as \emph{passive
followers} are monostable at rest and self-reset; cells designated as \emph{intrinsic oscillators}
are bistable and require network-mediated return.

% ──────────────────────────────────────────────────────────────────────────────
\section{Network Connectivity Analysis}
\label{sec:connectivity}

The 102-cell network contains nine cell classes.  Classes 1--7 are interneurons and motor neurons;
classes 8--9 are muscle outputs.  The class identities are: class~1 AS, class~2 DA, class~3 DB
(dorsal excitatory cholinergic); class~4 DD, class~5 VD (GABAergic inhibitory); class~6 VB,
class~7 VA (ventral excitatory cholinergic); classes 8--9 muscle cells.  The class-level
connectivity extracted from the six-segment connectivity matrices is summarised in
Table~\ref{tab:connectivity}.

\begin{table}[ht]
\centering
\caption{Class-level synaptic and gap-junction connectivity.
Entries show target class and total connection weight (sum over all pre/post pairs).}
\label{tab:connectivity}
\renewcommand{\arraystretch}{1.3}
\begin{tabular}{ccrllll}
\toprule
Class & $N$ & \multicolumn{1}{c}{Cell ID} & Exc.\ output & Inh.\ output & GJ partners & Inh.\ input \\
\midrule
1 & 12 & AS     & 2,4,5,8 & --- & 2, 7 & \textbf{none} \\
2 &  6 & DA     & 3,4,5,8 & --- & 1, 7 & \textbf{none} \\
3 &  6 & DB     & 1,4,5,8 & --- & 3, 6 & \textbf{none} \\
4 &  6 & DD     & ---      & 5, 8 & 4, 5 & \textbf{none} \\
5 & 12 & VD     & ---      & 6,7,9 & 4,5,7 & class~4 \\
6 & 12 & VB     & 4,5,7,9 & --- & 3, 6 & class~5 \\
7 & 12 & VA     & 4,5,6,9 & --- & 1,2,5 & class~5 \\
8 & 18 & muscle & ---      & --- & --- & class~4 \\
9 & 18 & muscle & ---      & --- & --- & class~5 \\
\bottomrule
\end{tabular}
\end{table}

The two inhibitory classes are:
\begin{itemize}
    \item \textbf{Class~4 (DD, 6 cells)}: inhibits class~5 (VD) and muscle group~8.
          Receives \emph{no} inhibitory input; can only return to rest if it is monostable.
    \item \textbf{Class~5 (VD, 12 cells)}: inhibits classes~6, 7 (VB, VA motor neurons) and muscle
          group~9.  Receives weak inhibitory input from class~4.
\end{itemize}

% ──────────────────────────────────────────────────────────────────────────────
\section{The Return-to-Rest Problem}
\label{sec:return}

\subsection{Two excitatory groups with asymmetric inhibitory coverage}

The excitatory motor neurons split into two groups based on which muscles they drive and which
inhibitory class they interact with:

\begin{itemize}
    \item \textbf{Group~A} (classes 1, 2, 3 = AS, DA, DB --- dorsal excitatory): drive muscle
          group~8; excite inhibitory class~4 (DD).  Receive \emph{no direct inhibitory synaptic
          input}.
    \item \textbf{Group~B} (classes 6, 7 = VB, VA --- ventral excitatory): drive muscle group~9;
          excite inhibitory class~5 (VD).  Receive \emph{direct inhibitory input from class~5
          (VD)}.
\end{itemize}

This asymmetry follows directly from the connectivity matrices.  VD (class~5) makes inhibitory
synapses onto the ventral motor neurons VB, VA (classes~6, 7) and onto the ventral muscle
(group~9).  DD (class~4) makes inhibitory synapses onto VD (class~5) and onto the dorsal muscle
(group~8), but \emph{not} onto the dorsal motor neurons AS, DA, DB.  Consequently the dorsal
excitatory group (classes~1, 2, 3) has no direct inhibitory synapse onto it.  Note that each
inhibitory neuron targets the muscle on the same side as the excitatory group that drives it
(DD$\to$dorsal muscle, VD$\to$ventral muscle), unlike the cross-inhibition seen in the biological
\textit{C.~elegans} circuit.

\subsection{Indirect return paths for Group~A}

With no direct inhibitory input, Group~A neurons can only be returned to rest via gap-junction
coupling.  The available paths are:

\begin{align*}
    \text{Class~1 (AS)} &\xrightarrow{\mathrm{GJ}} \text{Class~7 (VA)} \xleftarrow{\mathrm{inh}} \text{Class~5 (VD)} \\
    \text{Class~2 (DA)} &\xrightarrow{\mathrm{GJ}} \text{Class~7 (VA)} \xleftarrow{\mathrm{inh}} \text{Class~5 (VD)} \\
    \text{Class~3 (DB)} &\xrightarrow{\mathrm{GJ}} \text{Class~6 (VB)} \xleftarrow{\mathrm{inh}} \text{Class~5 (VD)}
\end{align*}

For this mechanism to work, VD must inhibit class~6 or~7 strongly enough to drive their voltage
below $T$, and the gap-junction current must then pull classes~1, 2, or~3 below $T$ as well.
The gap-junction current from cell $j$ onto cell $i$ is
\begin{equation}
    I_{\mathrm{GJ},i} = G_{\mathrm{gap}}\,(V_j - V_i).
\end{equation}
For this current to pull cell~$i$ below $T$ while it sits near $H$, the partner cell~$j$ must
be well below $T$, and $G_{\mathrm{gap}}$ must be large enough to overcome the restoring force of
$g\,f(V_i)$ near the plateau.

\subsection{Class~3 is the weakest link}

Class~3 (DB) is gap-junction coupled only to itself (intra-class, weight 10) and to class~6 (VB,
weight 5).  The GJ weight from class~6 to class~3 is the smallest single inter-class GJ weight in
the network.  If class~6 (VB) is successfully returned to rest by VD, the coupling onto class~3
(DB) must still be strong enough to drag it below $T$.  Class~3 (DB) is therefore the cell most
likely to remain latched at the plateau when the rest of the network returns to rest.

% ──────────────────────────────────────────────────────────────────────────────
\section{Parameter Constraints for Oscillation}

The analysis above implies a narrow feasibility corridor in the conductance parameter space
$(G_{\mathrm{syne}},\, G_{\mathrm{syni}},\, G_{\mathrm{gap}})$.

\paragraph{Lower bound on $G_{\mathrm{syni}}$.}
VD (class~5) must inhibit classes~6 and~7 strongly enough to drive them below $T$.  The required
inhibitory current at the plateau is approximately
\begin{equation}
    G_{\mathrm{syni}}\,\sigma(V_{\mathrm{pre}})\,(H - E_{\mathrm{syni}})
    \gtrsim g\,f(V)\big|_{V \to T^+} + I_{\mathrm{app}},
\end{equation}
where $\sigma$ is the presynaptic sigmoid gating variable evaluated at the active VD voltage.

\paragraph{Lower bound on $G_{\mathrm{gap}}$.}
Once classes~6 and~7 are returned to rest, the gap-junction current must be sufficient to pull
classes~1, 2, and~3 below $T$:
\begin{equation}
    G_{\mathrm{gap}}\,(V_{\mathrm{partner}} - H) \gtrsim -g\,f(H) - I_{\mathrm{app}}.
\end{equation}
Since $f(H) = 0$, the right-hand side reduces to $-I_{\mathrm{app}}$, so larger $I_{\mathrm{app}}$
demands proportionally larger $G_{\mathrm{gap}}$.

\paragraph{Upper bound on $G_{\mathrm{gap}}$.}
Gap junctions also couple cells within the same segment and between adjacent segments.  If
$G_{\mathrm{gap}}$ is too large, the coupling synchronises activity across all segments,
eliminating the anterior-to-posterior phase delay required for forward locomotion.

\paragraph{Summary.}
The feasibility conditions can be written schematically as:
\begin{equation}
    G_{\mathrm{syni}} > G_{\mathrm{syni}}^{\min}(\sigma, I_{\mathrm{app}}),
    \qquad
    G_{\mathrm{gap}}^{\min}(I_{\mathrm{app}}) < G_{\mathrm{gap}} < G_{\mathrm{gap}}^{\max}(\Delta\phi),
\end{equation}
where $\Delta\phi$ is the target inter-segment phase difference.  These bounds are simultaneously
satisfiable only in a narrow region of parameter space, consistent with the difficulty observed in
numerical sweeps.

% ──────────────────────────────────────────────────────────────────────────────
\section{Open Questions}

\begin{enumerate}
    \item \textbf{OscillatorComb selection}: which subset of classes 1--7 should be designated as
          intrinsic oscillators?  Designating Group~A cells (1, 2, 3) as passive followers (low
          $I_{\mathrm{app}}$) would remove the need for GJ-mediated return, at the cost of making
          them purely driven.
    \item \textbf{Class~4 (DD) return}: DD receives no inhibitory input.  If it is bistable, it
          will latch and continuously suppress VD (class~5), preventing oscillation.  It likely
          needs to be a passive follower.
    \item \textbf{Boyle vs.\ Yuval synapses}: the steep sigmoid of the Boyle config makes
          inhibitory gating more binary, potentially strengthening the return kick at lower
          $G_{\mathrm{syni}}$ values.  This may widen the feasibility corridor.
    \item \textbf{Sweep strategy}: given the narrow corridor, an informed sweep should fix
          $G_{\mathrm{syni}}$ above its estimated lower bound and sweep $(G_{\mathrm{gap}},
          G_{\mathrm{syne}})$, rather than treating all three as free simultaneously.
\end{enumerate}

\end{document}
